\setcounter{chapter}{9}\mxunit{terminale-c02}{Vecteurs, droites et plans de l'espace}\label{ch:10}

{\small\sffamily\color{mxmint}S3, S6, S12, S15 ; synthèse S17\par}

\begin{keypoint}[Objectifs]Calculer avec des vecteurs ; utiliser colinéarité et coplanarité ; caractériser droites et plans ; travailler dans un repère.\end{keypoint}

\section{Notes de cours}
\subsection{Calcul vectoriel}
La relation de Chasles $\vect{AB}+\vect{BC}=\vect{AC}$, la somme de vecteurs et le produit par un réel se prolongent à l'espace. Deux vecteurs non nuls sont colinéaires si l'un est un multiple de l'autre. Le vecteur nul est colinéaire à tout vecteur.
Trois vecteurs sont coplanaires s'ils appartiennent à un même plan vectoriel. Si $\vec u,\vec v$ ne sont pas colinéaires, $\vec u,\vec v,\vec w$ sont coplanaires si et seulement si $\vec w=\alpha\vec u+\beta\vec v$ pour certains réels $\alpha,\beta$.
\begin{keypoint}[Droite et plan]
La droite passant par $A$ de vecteur directeur non nul $\vec u$ est l'ensemble des points $M$ tels que $\vect{AM}=t\vec u$ pour un réel $t$.
Le plan passant par $A$ et dirigé par deux vecteurs non colinéaires $\vec u,\vec v$ est l'ensemble des points $M$ tels que $\vect{AM}=s\vec u+t\vec v$ pour des réels $s,t$.
\end{keypoint}
Trois points non alignés définissent un plan. Quatre points $A,B,C,D$, avec $A,B,C$ non alignés, sont coplanaires si et seulement si $\vect{AD}$ est combinaison linéaire de $\vect{AB}$ et $\vect{AC}$.
\subsection{Positions relatives}
Deux droites sont confondues, strictement parallèles, sécantes ou non coplanaires. Des vecteurs directeurs colinéaires caractérisent des droites parallèles au sens large. Des directeurs non colinéaires n'assurent pas l'intersection.
Une droite et un plan : droite incluse dans le plan, parallèle et disjointe, ou sécante en un seul point. Deux plans : confondus, strictement parallèles, ou sécants suivant une droite.
\paragraph{Exemple dans un cube.} Dans $ABCDEFGH$, les faces $ABCD$ et $EFGH$ sont parallèles et les arêtes $AE,BF,CG,DH$ se correspondent. Les droites $(AB)$ et $(CG)$ ne sont ni parallèles ni sécantes : elles sont non coplanaires.
\subsection{Base et coordonnées}
Une base de l'espace est une famille de trois vecteurs non coplanaires. Tout vecteur s'y écrit de manière unique $x\vi +y\vj +z\vec k$. Un repère est constitué d'une origine $O$ et d'une base.
\[\vect{AB}=(x_B-x_A,y_B-y_A,z_B-z_A),\qquad
I=\left(\frac{x_A+x_B}2,\frac{y_A+y_B}2,\frac{z_A+z_B}2\right).\]
Dans un repère orthonormé, $\|\vec u\|=\sqrt{x^2+y^2+z^2}$ et $AB=\sqrt{(x_B-x_A)^2+(y_B-y_A)^2+(z_B-z_A)^2}$.
\paragraph{Exemple.} Dans un repère orthonormé, pour $A(1,0,2)$, $B(3,-1,4)$, $\vect{AB}=(2,-1,2)$, $AB=3$ et le milieu est $(2,-\frac{1}{2},3)$.
\begin{method}[Tester une combinaison linéaire]
Écrire les trois égalités de coordonnées. Déterminer les coefficients à partir de deux équations quand c'est possible, puis vérifier la troisième. Une contradiction prouve la non-appartenance ; on ne supprime pas l'équation gênante.
\end{method}



\section{Activité de recherche}\textbf{Trois milieux dans un cube --- 40 min.}\par On considère le cube $ABCDEFGH$ représenté ci-dessous. On pose $\vi=\vect{AB}$, $\vj=\vect{AD}$ et $\vec k=\vect{AE}$. Les points $I$, $J$ et $K$ sont les milieux respectifs de $[AB]$, $[AD]$ et $[AE]$. On note $O$ le centre du cube.
\begin{center}\begin{tikzpicture}[scale=1.1,line join=round]
\coordinate(A)at(0,0);\coordinate(B)at(4,0);\coordinate(C)at(5.1,1.6);\coordinate(D)at(1.1,1.6);\coordinate(E)at(0,2.5);\coordinate(F)at(4,2.5);\coordinate(G)at(5.1,4.1);\coordinate(H)at(1.1,4.1);
\draw[mxink,thick](A)--(B)--(C)--(G)--(H)--(E)--(A) (E)--(F)--(G) (B)--(F);
\draw[mxink,dashed](A)--(D)--(C) (D)--(H);
\draw[mxmint,thick,dash pattern=on 5pt off 2pt](A)--(G);
\foreach \p/\pos in {A/below left,B/below,C/right,D/left,E/left,F/right,G/right,H/above}{\node[\pos]at(\p){$\p$};}
\end{tikzpicture}\par\small $[AG]$ est une diagonale de l'espace, tracée en pointillés verts.\\Les arêtes cachées sont en pointillés bleus ; $D$ et $F$ ne sont pas sur $(AG)$.\end{center}
\begin{enumerate}\item Exprimer $\vect{IJ}$ et $\vect{IK}$ dans la base $(\vi,\vj,\vec k)$.
\item Démontrer que les points $I$, $J$ et $K$ ne sont pas alignés.
\item Déterminer si $O$ appartient au plan $(IJK)$. On justifiera la réponse en cherchant si $\vect{IO}$ est une combinaison linéaire de $\vect{IJ}$ et $\vect{IK}$.
\item Pourquoi la seule lecture de la figure ne permet-elle pas de démontrer l'appartenance ou la non-appartenance de $O$ au plan $(IJK)$ ?\end{enumerate}

\section{Exercices progressifs}

\mxexerciseitem{1}{Entraînement}Dans une base de l'espace, les vecteurs $\vec u$ et $\vec v$ ont pour coordonnées respectives $(1;2;-1)$ et $(3;6;-3)$. Sont-ils colinéaires ? Justifier.

\mxexerciseitem{2}{Entraînement}Dans un repère orthonormé de l'espace, on considère $A(0;1;2)$ et $B(4;-1;0)$. Déterminer les coordonnées du milieu de $[AB]$, puis calculer la distance $AB$.

\mxexerciseitem{3}{Entraînement}Dans une base de l'espace, on considère $\vec u=(1;0;1)$, $\vec v=(0;1;1)$ et $\vec w=(2;3;5)$. Exprimer $\vec w$ comme combinaison linéaire de $\vec u$ et $\vec v$, puis en déduire que ces trois vecteurs sont coplanaires.

\mxexerciseitem{4}{Entraînement}Dans une base de l'espace, les vecteurs $\vec u$, $\vec v$ et $\vec w$ ont pour coordonnées $(1;0;1)$, $(0;1;1)$ et $(1;1;0)$. Le triplet $(\vec u,\vec v,\vec w)$ est-il une base de l'espace ? Justifier.

\mxexerciseitem{5}{Entraînement}Dans un repère de l'espace, on considère $A(1;0;2)$, $B(3;-1;4)$ et $C(5;-2;6)$. Déterminer si $C$ appartient à la droite $(AB)$. Justifier par un calcul vectoriel.

\section{Travail dans le manuel ou une annale}\textbf{Barbazo, p. 328, exercice 21.} En classe après S15 : traiter les trois questions. Pour chaque intersection, commencer par identifier chacun des plans à partir des sommets du cube et des milieux indiqués. Distinguer deux plans sécants, confondus et strictement parallèles.

\section{Devoir à la maison}\textbf{DM10 --- Points d'un tétraèdre --- 50 min.}\par On considère un tétraèdre $ABCD$ non aplati. Les points $I$, $J$ et $K$ sont les milieux respectifs de $[AB]$, $[AC]$ et $[AD]$.
\begin{enumerate}\item Exprimer $\vect{IJ}$ en fonction de $\vect{BC}$ et $\vect{IK}$ en fonction de $\vect{BD}$.
\item En déduire que les plans $(IJK)$ et $(BCD)$ sont parallèles.
\item On définit le point $G$ par $\vect{AG}=\frac14(\vect{AB}+\vect{AC}+\vect{AD})$. Le point $G$ appartient-il au plan $(IJK)$ ? Justifier.
\item Pour cette question seulement, on suppose que, dans un repère orthonormé, $A(0;0;0)$, $B(4;0;0)$, $C(0;4;0)$ et $D(0;0;4)$. Calculer les longueurs $IJ$, $IK$ et $JK$.\end{enumerate}

\ifmxteacher\mxnewpage\section{Conduite et corrigés}\begin{teachernote}[Repères pour le professeur]Le DS1 précède S12 et S15 : proposer de la géométrie vectorielle dans une base donnée, sans exiger de techniques de coordonnées nouvelles. Utiliser un cube matériel, puis une perspective cavalière ; marquer les arêtes cachées. Deux droites non parallèles de l'espace ne sont pas nécessairement sécantes. La norme par coordonnées nécessite un repère orthonormé. Les représentations paramétriques formalisées et les équations de plans seront traitées plus tard.\end{teachernote}

\subsection{Corrigé de l'activité}\begin{enumerate}
\item Par Chasles, $\vect{IJ}=\vect{AJ}-\vect{AI}=-\frac12\vi +\frac12\vj $ et $\vect{IK}=-\frac12\vi +\frac12\vec k$.
\item Si $\vect{IK}=t\vect{IJ}$, le coefficient de $\vec k$ serait $\frac{1}{2}=0$, impossible dans une base. Ces deux vecteurs sont donc non colinéaires ; $I,J,K$ définissent un plan.
\item Le centre du cube vérifie $\vect{AO}=\frac12(\vi +\vj +\vec k)$, donc $\vect{IO}=\frac12\vj +\frac12\vec k$. Si $O\in(IJK)$, il existe $s,t$ tels que $\vect{IO}=s\vect{IJ}+t\vect{IK}$. En comparant les coefficients de la base,
\[0=-\frac{s+t}2,\qquad \frac12=\frac s2,\qquad \frac12=\frac t2.\]
Les deux dernières imposent $s=t=1$, alors que la première demanderait $s+t=0$. Contradiction : $O\notin(IJK)$.
\item La perspective cavalière est une projection sur une feuille : des points distincts peuvent avoir la même image. Elle aide à conjecturer, mais une superposition apparente ne prouve pas l'appartenance à un plan.
\end{enumerate}

\paragraph{E1.}Les trois coordonnées de $\vec v$ valent trois fois celles de $\vec u$ : $\vec v=3\vec u$. Les deux vecteurs sont donc colinéaires.

\paragraph{E2.}Le milieu a pour coordonnées $(\frac{0+4}{2},\frac{1-1}{2},\frac{2+0}{2})=(2,0,1)$. La différence de coordonnées donne $\vect{AB}=(4,-2,-2)$ ; dans le repère orthonormé, $AB=\sqrt{16+4+4}=2\sqrt6$.

\paragraph{E3.}On calcule $2\vec u+3\vec v=(2,0,2)+(0,3,3)=(2,3,5)=\vec w$. Ainsi $\vec w$ appartient au plan vectoriel dirigé par les deux premiers vecteurs : les trois sont coplanaires.

\paragraph{E4.}Les deux premiers vecteurs sont non colinéaires. Si le troisième était leur combinaison $\lambda(1,0,1)+\mu(0,1,1)$, les deux premières coordonnées imposeraient $\lambda=\mu=1$, mais la troisième vaudrait $2\ne0$. Les trois vecteurs sont non coplanaires ; ils forment donc une base de l'espace.

\paragraph{E5.}$\vect{AB}=(2,-1,2)$ et $\vect{AC}=(4,-2,4)=2\vect{AB}$. Comme $A\ne B$, la droite $(AB)$ est définie et cette colinéarité prouve $C\in(AB)$.

\subsection{Corrigé du DM10}\begin{enumerate}
\item Par Chasles et les définitions des milieux,
\[\vect{IJ}=\tfrac12\vect{AC}-\tfrac12\vect{AB}=\tfrac12\vect{BC},\qquad \vect{IK}=\tfrac12\vect{AD}-\tfrac12\vect{AB}=\tfrac12\vect{BD}.\]
\item Les points $B,C,D$ ne sont pas alignés, car le tétraèdre est non aplati. Les vecteurs $\vect{BC},\vect{BD}$ sont donc non colinéaires, de même que $\vect{IJ},\vect{IK}$. Les plans ont la même direction. Ils sont distincts : si $I\in(BCD)$, la droite $(BI)=(BA)$ serait incluse dans ce plan, ce qui y placerait $A$, contradiction. Ils sont donc strictement parallèles.
\item Dans la base $(\vect{AB},\vect{AC},\vect{AD})$, on a $\vect{IG}=(-\frac{1}{4},\frac{1}{4},\frac{1}{4})$, $\vect{IJ}=(-\frac{1}{2},\frac{1}{2},0)$ et $\vect{IK}=(-\frac{1}{2},0,\frac{1}{2})$. Une relation $\vect{IG}=s\vect{IJ}+t\vect{IK}$ impose $s=t=\frac{1}{2}$ par les deux dernières coordonnées ; la première vaudrait alors $-\frac{1}{2}$, et non $-\frac{1}{4}$. Donc $G\notin(IJK)$.
\item Dans ce cas particulier, $I(2,0,0)$, $J(0,2,0)$ et $K(0,0,2)$. On trouve $IJ=\sqrt{(-2)^2+2^2}=2\sqrt2$, et de même $IK=JK=2\sqrt2$. Le triangle est équilatéral. Cette conclusion métrique ne vaut pas pour tout tétraèdre non aplati : elle utilise les coordonnées du cas particulier.
\end{enumerate}

\subsection{Corrigé du travail dans le manuel ou une annale}Dans cet exercice, $I$ est le milieu de $[AB]$ et $J$ celui de $[CD]$.
\begin{enumerate}
\item Les plans $(AIE)$ et $(BIG)$ contiennent tous deux $(AB)$. Ils sont distincts : $G$ n'appartient pas au plan de la face $ABFE$. Leur intersection est donc exactement la droite $(AB)$.
\item Les points $A,D,I$ sont trois points non alignés de la face $ABCD$ ; il en va de même de $B,J,C$. Les deux plans sont donc confondus avec $(ABC)$. L'intersection est le plan entier, et non une droite.
\item Le plan $(HEF)$ est celui de la face supérieure ; $(BJC)$ est celui de la face inférieure. Ces plans sont strictement parallèles : leur intersection est vide.
\end{enumerate}

\fi

